\documentclass[10pt,a4paper]{amsart}
\usepackage{amsmath,amssymb,amsthm,graphicx}
% Retained minimal preamble for the Stacks Project chapter more-algebra.tex (More on Algebra).
% The chapter file pinned-original-source.tex carries only \input{preamble} on its first line, so no
% preamble text lives inside the chapter itself.  The Stacks Project uses one single shared preamble for
% every chapter; the snapshot stages it identically as sources/stacks-chow/preamble.tex and
% sources/stacks-etale-cohomology/preamble.tex.  The block below is copied mechanically, line for line,
% from sources/stacks-chow/preamble.tex lines 14-18 (the Xy-pic setup) and 154-217 (the native theorem
% declarations together with the mathematical macros).  Nothing was added, removed or redefined; the
% wrapper class and the amsmath/amssymb/amsthm/graphicx base packages are supplied by the assembler.
% Chapter: more-algebra.tex -- More on Algebra.
% For commutative diagrams we use Xy-pic
\usepackage[all]{xy}

% We use 2cell for 2-commutative diagrams.
\xyoption{2cell}

% Theorem environments.
%
\theoremstyle{plain}
\newtheorem{theorem}[subsection]{Theorem}
\newtheorem{proposition}[subsection]{Proposition}
\newtheorem{lemma}[subsection]{Lemma}

\theoremstyle{definition}
\newtheorem{definition}[subsection]{Definition}
\newtheorem{example}[subsection]{Example}
\newtheorem{exercise}[subsection]{Exercise}
\newtheorem{situation}[subsection]{Situation}

\theoremstyle{remark}
\newtheorem{remark}[subsection]{Remark}
\newtheorem{remarks}[subsection]{Remarks}

\numberwithin{equation}{subsection}

% Macros
%
\def\lim{\mathop{\mathrm{lim}}\nolimits}
\def\colim{\mathop{\mathrm{colim}}\nolimits}
\def\Spec{\mathop{\mathrm{Spec}}}
\def\Hom{\mathop{\mathrm{Hom}}\nolimits}
\def\Ext{\mathop{\mathrm{Ext}}\nolimits}
\def\SheafHom{\mathop{\mathcal{H}\!\mathit{om}}\nolimits}
\def\SheafExt{\mathop{\mathcal{E}\!\mathit{xt}}\nolimits}
\def\Sch{\mathit{Sch}}
\def\Mor{\mathop{\mathrm{Mor}}\nolimits}
\def\Ob{\mathop{\mathrm{Ob}}\nolimits}
\def\Sh{\mathop{\mathit{Sh}}\nolimits}
\def\NL{\mathop{N\!L}\nolimits}
\def\CH{\mathop{\mathrm{CH}}\nolimits}
\def\proetale{{pro\text{-}\acute{e}tale}}
\def\etale{{\acute{e}tale}}
\def\QCoh{\mathit{QCoh}}
\def\Ker{\mathop{\mathrm{Ker}}}
\def\Im{\mathop{\mathrm{Im}}}
\def\Coker{\mathop{\mathrm{Coker}}}
\def\Coim{\mathop{\mathrm{Coim}}}

% Boxtimes
%
\DeclareMathSymbol{\boxtimes}{\mathbin}{AMSa}{"02}

%
% Macros for moduli stacks/spaces
%
\def\QCohstack{\mathcal{QC}\!\mathit{oh}}
\def\Cohstack{\mathcal{C}\!\mathit{oh}}
\def\Spacesstack{\mathcal{S}\!\mathit{paces}}
\def\Quotfunctor{\mathrm{Quot}}
\def\Hilbfunctor{\mathrm{Hilb}}
\def\Curvesstack{\mathcal{C}\!\mathit{urves}}
\def\Polarizedstack{\mathcal{P}\!\mathit{olarized}}
\def\Complexesstack{\mathcal{C}\!\mathit{omplexes}}
% \Pic is the operator that assigns to X its picard group, usage \Pic(X)
% \Picardstack_{X/B} denotes the Picard stack of X over B
% \Picardfunctor_{X/B} denotes the Picard functor of X over B
\def\Pic{\mathop{\mathrm{Pic}}\nolimits}
\def\Picardstack{\mathcal{P}\!\mathit{ic}}
\def\Picardfunctor{\mathrm{Pic}}
\def\Deformationcategory{\mathcal{D}\!\mathit{ef}}

% --- appended 2026-10-12 (Chief): the project own definitions of the reference/slogan/history
% display environments, copied mechanically from the shared Stacks preamble lines 8-12
% (\usepackage{verbatim} plus the three \newenvironment lines, which the chapter preamble keeps
% verbatim from the project preamble). Nothing added, removed or redefined: these environments
% carry no mathematical content and the project itself typesets them as comments. Without these
% three lines a counted statement whose own text uses them fails with Environment slogan undefined.
\usepackage{verbatim}
\newenvironment{reference}{\comment}{\endcomment}
\newenvironment{slogan}{\comment}{\endcomment}
\newenvironment{history}{\comment}{\endcomment}
\usepackage{cleveref}
\crefname{originallemma}{original lemma}{original lemmas}
\Crefname{originallemma}{Original lemma}{Original lemmas}
\makeatletter
\newlabel{more-algebra-lemma-lci-local}{{07D1}{source}{}{sourcepointer-1}{}}
\newlabel{more-algebra-lemma-lci-local@cref}{{[originallemma][0][]07D1}{[1][0][]source}{}{}{}}
\makeatother
\crefname{originallemma}{original lemma}{original lemmas}
\Crefname{originallemma}{Original lemma}{Original lemmas}
\makeatletter
\newlabel{algebra-lemma-principal-localization-NL}{{08JZ}{source}{}{sourcepointer-2}{}}
\newlabel{algebra-lemma-principal-localization-NL@cref}{{[originallemma][0][]08JZ}{[1][0][]source}{}{}{}}
\makeatother
\crefname{originallemma}{original lemma}{original lemmas}
\Crefname{originallemma}{Original lemma}{Original lemmas}
\makeatletter
\newlabel{lemma-independence-of-generators}{{066A}{source}{}{sourcepointer-3}{}}
\newlabel{lemma-independence-of-generators@cref}{{[originallemma][0][]066A}{[1][0][]source}{}{}{}}
\makeatother
\crefname{originallemma}{original lemma}{original lemmas}
\Crefname{originallemma}{Original lemma}{Original lemmas}
\makeatletter
\newlabel{lemma-truncate-koszul-regular}{{068M}{source}{}{sourcepointer-4}{}}
\newlabel{lemma-truncate-koszul-regular@cref}{{[originallemma][0][]068M}{[1][0][]source}{}{}{}}
\makeatother
\title{The local complete intersection property for a ring map is local on the target in the Zariski topology}
\author{The Stacks Project Authors}
\date{stacks-project:more-algebra.tex:local-snapshot-20261012; GNU FDL 1.2 or later}
\begin{document}
\maketitle
\setcounter{subsection}{1}
\noindent\textit{Editorial note.} All delivered mathematical text is the author's own, copied from the pinned Stacks Project chapter file more-algebra.tex; no statement or proof step was rewritten, reordered or removed.\par\medskip
\noindent\textit{Editorial note.} The chapter file pinned-original-source.tex carries only \textbackslash{}input\{preamble\} on its first line, so no preamble text lives inside the chapter. The retained file sources/stacks-more-algebra/retained-preamble.tex is a mechanical, line-for-line copy of the project's single shared preamble as staged in the snapshot (Xy-pic setup plus theorem declarations and mathematical macros). Nothing was added, removed or redefined, and the excerpt serves only as preamble, never as a span coordinate.\par\medskip
\noindent\textit{Editorial note.} Both counted spans start at the \textbackslash{}begin line of their environment, so the statement and the author's proof are delivered whole; surrounding prose of the section is deliberately outside the counted spans and is carried as standard prerequisites instead.\par\medskip
\noindent\textit{Editorial note.} The citation fields (source\_id, source\_version, authors, source\_url, license.name, source\_title) are identical to those of every other card this lane has produced from this chapter, so the batch selector treats them as one source.\par\medskip
\noindent\textit{Allowed standard prerequisites (editorial).}\begin{itemize}
\item Local complete intersection for a ring map, expressed through the Koszul regularity of the ideal presenting a localisation. used throughout the proof, source lines 8354-8383. Supplies the hypothesis at each localisation and the goal at the original ring.
\item Independence of the choice of generators for a Koszul-regular ideal. the statement counted by the sibling card of this lane for Lemma lemma-independence-of-generators at 7711-7755, cited at proof lines 8375-8378. Turns the generating Koszul-regular sequence into the specific sequence obtained from the conormal basis.
\item Truncation of a Koszul-regular sequence into a quotient. the statement counted by the sibling card S1-algebra-truncate-koszul-regular at 7679-7709, cited at proof lines 8378-8381. Descends Koszul regularity from the extended presentation to the original one.
\end{itemize}\medskip
\ifdefined\hypertarget\hypertarget{sourcepointer-1}{}\fi
\noindent\textit{Source reference.} Original lemma 07D1: focused-seeds/AGP/sources/stacks-snapshot-20261011/tags.txt, source lines 9396--9396. The lemma's own canonical tag, quoted so that readers can resolve the label inside the delivered text; it points at the counted statement itself and carries no proof content.\par\medskip
\ifdefined\hypertarget\hypertarget{sourcepointer-2}{}\fi
\noindent\textit{Source reference.} Original lemma 08JZ: focused-seeds/AGP/sources/stacks-snapshot-20261011/tags.txt, source lines 10867--10867. Navigation pointer to a statement the counted proof refers to but does not contain; nothing load-bearing is delegated to it.\par\medskip
\ifdefined\hypertarget\hypertarget{sourcepointer-3}{}\fi
\noindent\textit{Source reference.} Original lemma 066A: focused-seeds/AGP/sources/stacks-snapshot-20261011/tags.txt, source lines 7929--7929. Navigation pointer to a statement the counted proof refers to but does not contain; nothing load-bearing is delegated to it.\par\medskip
\ifdefined\hypertarget\hypertarget{sourcepointer-4}{}\fi
\noindent\textit{Source reference.} Original lemma 068M: focused-seeds/AGP/sources/stacks-snapshot-20261011/tags.txt, source lines 8011--8011. Navigation pointer to a statement the counted proof refers to but does not contain; nothing load-bearing is delegated to it.\par\medskip
\begin{lemma}
\label{lemma-lci-local}
Let $R \to S$ be a ring map. Let $g_1, \ldots, g_m \in S$
generate the unit ideal. If each $R \to S_{g_j}$ is a local
complete intersection so is $R \to S$.
\end{lemma}

\begin{proof}
Let $S = R[x_1, \ldots, x_n]/I$ be a presentation. Pick
$h_j \in R[x_1, \ldots, x_n]$ mapping to $g_j$ in $S$.
Then $R[x_1, \ldots, x_n, x_{n + 1}]/(I, x_{n + 1}h_j - 1)$
is a presentation of $S_{g_j}$. Hence
$I_j = (I, x_{n + 1}h_j - 1)$ is a Koszul-regular ideal in
$R[x_1, \ldots, x_n, x_{n + 1}]$. Pick a prime
$I \subset \mathfrak q \subset R[x_1, \ldots, x_n]$.
Then $h_j \not \in \mathfrak q$ for some $j$ and
$\mathfrak q_j = (\mathfrak q, x_{n + 1}h_j - 1)$ is a prime
ideal of $V(I_j)$ lying over $\mathfrak q$.
Pick $f_1, \ldots, f_r \in I$ which map to a basis of
$I/I^2 \otimes \kappa(\mathfrak q)$. Then
$x_{n + 1}h_j - 1, f_1, \ldots, f_r$ is a sequence of elements of $I_j$
which map to a basis of $I_j \otimes \kappa(\mathfrak q_j)$, see
Algebra, Lemma \ref{algebra-lemma-principal-localization-NL}.
By Nakayama's lemma there exists an $h \in R[x_1, \ldots, x_n, x_{n + 1}]$
such that $(I_j)_h$ is generated by $x_{n + 1}h_j - 1, f_1, \ldots, f_r$.
We may also assume that $(I_j)_h$ is generated by a Koszul regular
sequence of some length $e$. Looking at the dimension of
$I_j \otimes \kappa(\mathfrak q_j)$ we see that $e = r + 1$.
Hence by
Lemma \ref{lemma-independence-of-generators}
we see that $x_{n + 1}h_j - 1, f_1, \ldots, f_r$ is a
Koszul-regular sequence generating $(I_j)_h$ for some
$h \in R[x_1, \ldots, x_n, x_{n + 1}]$, $h \not \in \mathfrak q_j$. By
Lemma \ref{lemma-truncate-koszul-regular}
we see that $I_{h'}$ is generated by a Koszul-regular sequence for some
$h' \in R[x_1, \ldots, x_n]$, $h' \not \in \mathfrak q$ as desired.
\end{proof}

\end{document}
